\scanpage{256}
\mspage{206}

est une vertu particulière des \unsure{Vogons|Vosges|Vagos}, qui ne doit pas pour autant faire oublier le caractère intrinsèque des flèches \unsure{(x)|(*)}, indépendamment de \unsure{tout} mode de calcul.

\margin{\rule{0.35pt}{4\baselineskip}\hspace{2pt}\rule{0.35pt}{4\baselineskip}}%
La clef de la situation, et des questions de canonicité, me semble être dans le fait que \unsure{(xx)|(**)} est une équivalence de catégories -- et il faut avouer que, tout familier et quasiment banal qu'il soit finalement devenu pour moi, c'est là un fait tout \unsure{à|:} fait remarquable et nullement évident a priori. (Mais sans doute peut-on \unsure{le avoir|l'avoir} aussi par les \unsure{fourbis} généraux à la Quillen-Thomason, sans se battre avec un formalisme des chemins ...) Les ``isom can'' qu'on peut établir entre une ``théorie $C$'' et une ``théorie $C'$'' (via deux constructions différentes, p. ex, de catégories de chemins, ... ou via l'approche de Quillen-Thomason, sans chemins) doivent toujours revenir à \unsure{celui|cela|celui-ci}
% archive sheet number: 256
